\documentclass{tietreport}

% Import images
\usepackage{graphicx}

% Bibliography configuration
\usepackage[
  date=year,
  style=alphabetic,
  backend=biber,
  sorting=nty,
  maxnames=5,
  minnames=3,
  maxalphanames=4,
  minalphanames=3,
  backref=true,
  doi=false,
  isbn=false,
  url=false,
  eprint=false
]{biblatex}
\addbibresource{references.bib}

% Placeholder box for diagrams that are still to be drawn.
% Replace each \diagramslot{...} with \includegraphics once the
% figure files exist in report/figures/.
% #1 = placeholder text, #2 = caption, #3 = label key, #4 = height
\newcommand{\diagramslot}[4]{%
  \begin{figure}[h]
    \centering
    \fbox{%
      \begin{minipage}[c][#4][c]{0.88\textwidth}
        \centering\itshape\small #1
      \end{minipage}%
    }
    \caption{#2}
    \label{fig:#3}
  \end{figure}%
}

%% ==================================================
%% PROJECT DETAILS
%% ==================================================

\TitlePageHeader{%
  UCS503P - Software Engineering Laboratory%
}

\ProjectTitle{%
  CareBridge%
}

\ProjectSubTitle{%
  A Structured Nurse Handoff and Patient Status
  Management System%
}

\author{%
  \texttt{102XXXXXXX} Arnav Sharma \\
  \texttt{102XXXXXXX} Dhruv Dwivedi \\
  \texttt{102XXXXXXX} Pratham Nandwani \\
  \texttt{102XXXXXXX} Jalaj Poonia%
}

\TitlePageSubText{%
  Group: \texttt{XXXX} BE Third Year, CSE%
}

\AdvisorName{%
  Anushka Choudhary%
}

\TitlePageFooterText{{%
    \large Mid-Semester Evaluation \\[0.2cm]
    \bfseries Computer Science and Engineering
    Department%
  } \\[0.15cm]
  Thapar Institute of Engineering and Technology,
  Patiala%
}

%% ==================================================

\begin{document}

\maketitle

\tableofcontents

% ===== CHAPTER 1: INTRODUCTION =====

\chapter{Introduction}

\section{Background and Context}

A nursing shift handoff is the moment responsibility for a
patient passes from one nurse to another. It happens several
times a day in every ward, and it is one of the few points
in clinical care where the full state of a patient exists
only in someone's short-term memory. What the outgoing nurse
remembers to say is what the incoming nurse knows.

This is a documented weak point rather than a hypothetical
one. The World Health Organization treats communication
during hand-overs as a priority patient-safety
problem~\cite{who2007handover}; reviews of the nursing
handoff literature find the process unstandardised and
variable between individuals~\cite{riesenberg2010handoffs};
and accreditation bodies have issued specific guidance on
inadequate hand-off
communication~\cite{jointcommission2017handoff}. Where
structured protocols have been introduced and studied, the
measured effect is a substantial reduction in preventable
adverse events~\cite{starmer2014ipass}.

CareBridge gives nurses a structured place to record patient
status through the shift, and an accountable mechanism for
transferring that record to the next nurse.

\subsection{Problem Statement}

Current ward practice shows five recurring failures:

\begin{enumerate}
  \item \textbf{Nothing is recorded consistently.} Handoffs
    happen verbally or from memory, so completeness depends
    on recall.
  \item \textbf{Quality varies with the individual.} The
    process has no floor: a rushed handoff at the end of a
    long shift may omit the detail that mattered.
  \item \textbf{There is no record of who said what.} Later,
    nothing establishes what was communicated, by whom, or
    whether it was acknowledged.
  \item \textbf{There is no way to push back.} An unclear or
    misdirected handoff cannot be formally rejected or
    flagged.
  \item \textbf{Nothing is measurable.} The handoff leaves
    no trace, so there is no way to tell whether it is
    improving.
\end{enumerate}

The common cause is that the handoff has no \emph{artefact}.
Nothing written down means nothing to review, audit, dispute
or measure. CareBridge creates that artefact and builds
accountability around it.

\section{Project Objectives}

\begin{enumerate}
  \item \textbf{Structured patient record.} Let a nurse
    accumulate notes per patient through the shift, split
    into \emph{Details} and \emph{Major Decisions}, so the
    handoff is composed from a record rather than from
    memory.
  \item \textbf{Per-patient, addressed handoff.} Let her
    choose which patients to hand over, edit the text for
    each, and send each to a \emph{named} incoming nurse.
  \item \textbf{Explicit accountability.} Require a
    response of \emph{Accept}, \emph{Decline} or
    \emph{Report Issue} --- with a mandatory reason for the
    latter two --- and record every transition in an
    append-only audit trail.
  \item \textbf{Measurability.} Make Time-to-Acknowledgement
    (TTA), the interval between sending and answering,
    computable directly from stored timestamps.
  \item \textbf{Verified access control.} Ensure a nurse can
    reach only her own patients, or one being handed to her,
    and demonstrate it with automated tests.
\end{enumerate}

\subsection{Scope}
\label{sec:scope}

The system defines three actors: nurse, doctor and
administrator. This deliverable implements the \textbf{nurse
workflow completely} and defers the other two, on the
principle that one working actor is demonstrable and
testable where three partial ones are neither. The schema
and access-control model already accommodate all three, so
the remaining interfaces are additive.
Section~\ref{sec:deferred} lists what is deferred.

% ===== CHAPTER 2: METHODOLOGY & DESIGN =====

\chapter{Methodology and Design}

\section{System Architecture}

CareBridge is a three-tier application over a managed
database.

\diagramslot{[DIAGRAM SLOT --- System architecture: client
tier (React/Tauri) $\rightarrow$ application tier (FastAPI)
$\rightarrow$ data tier (Supabase Postgres), with Supabase
Auth as an external dependency. Mark the two trust
boundaries: the client holds only the publishable key, the
server holds the secret key.]}{System
Architecture}{architecture}{6cm}

\begin{itemize}
  \item \textbf{Client:} a React single-page application in
    a Tauri shell. It holds no business rules.
  \item \textbf{Application:} a FastAPI service owning every
    rule --- who may read a patient, what a valid response
    is, when a shift may close. It is the only component
    that talks to the database.
  \item \textbf{Data:} Supabase-managed
    PostgreSQL~\cite{postgresql} with Row Level Security
    beneath the application's own checks.
\end{itemize}

The client never talks to the database directly, even though
Supabase allows it. Routing all access through FastAPI keeps
clinical rules in one auditable place and keeps the secret
key off client devices.

\subsection{Identity}

Identity comes from the Supabase Auth~\cite{supabase} JWT
and nothing else. The client signs in, and an interceptor
attaches the token to every request; the backend verifies it
and resolves it to a \texttt{staff} row.

No endpoint accepts a nurse identifier in a request body.
If one did, any caller could act as any nurse and the audit
trail would record only what the client chose to claim.

The frontend holds only the \emph{publishable} key, enough
to sign in. The backend holds the \emph{secret} key and so
bypasses RLS --- which is why it performs its own
authorisation, in \texttt{services/access.py}, applied as a
dependency on every route.

\subsection{Data Model}

\diagramslot{[DIAGRAM SLOT --- Entity-Relationship diagram
for the seven tables, showing cardinalities and the two
separate links from \texttt{patients} to \texttt{staff}
(assigned nurse, assigned doctor).]}{Entity Relationship
Diagram}{erd}{7cm}

\begin{table}[h]
\centering
{\small
\begin{tabular}{|l|p{8.2cm}|}
  \hline
  \textbf{Table} & \textbf{Purpose} \\
  \hline
  \texttt{staff} & All personnel, keyed to
    \texttt{auth.users}, with a nurse/doctor/admin role. \\
  \hline
  \texttt{patients} & Admitted patients, each with an
    assigned nurse and doctor. The assigned nurse is what
    makes ``my patients'' meaningful. \\
  \hline
  \texttt{shifts} & Shift lifecycle per nurse. \\
  \hline
  \texttt{patient\_notes} & The clinical record: body text,
    a section (\texttt{detail} or
    \texttt{major\_decision}), an author and a shift. \\
  \hline
  \texttt{handoffs} & One row per patient handed over ---
    sender, recipient, text, response, and the timestamps
    TTA derives from. \\
  \hline
  \texttt{handoff\_events} & Append-only audit trail:
    \texttt{sent}, \texttt{viewed}, \texttt{accepted},
    \texttt{declined}, \texttt{issue\_reported}. \\
  \hline
  \texttt{alerts} & Provisioned for nurse-to-doctor
    escalation, delivered with the doctor interface. \\
  \hline
\end{tabular}}
\caption{Database schema}
\label{tab:schema}
\end{table}

Two decisions here matter. \textbf{Notes are one table, not
two}: section is a column, so the same rows serve the
sectioned views and the chronological log and cannot
diverge. \textbf{A handoff is per patient, not per shift},
because accountability attaches to a patient --- which is
what makes it possible to decline one patient without
rejecting a whole batch.

\subsection{Rules Enforced by the Database}

Invariants that must never break are constraints, so they
hold regardless of which code path runs:

\begin{itemize}
  \item A handoff cannot be sent to oneself.
  \item Declining or reporting an issue requires a non-empty
    reason.
  \item A resolved handoff always has a
    \texttt{responded\_at}; a pending one never does ---
    keeping TTA well-defined by construction.
  \item At most one \emph{pending} handoff per patient, and
    one \emph{active} shift per nurse (partial unique
    indexes).
\end{itemize}

\subsection{Technical Stack}

\begin{itemize}
  \item \textbf{Frontend:} React 19~\cite{react},
    TypeScript, Vite, React Router, Tailwind CSS, in Tauri
    v2~\cite{tauri}.
  \item \textbf{Backend:} Python 3.13,
    FastAPI~\cite{fastapi}, Uvicorn, Pydantic v2.
  \item \textbf{Database and auth:} Supabase (PostgreSQL 17)
    with RLS; Supabase Auth for identity.
  \item \textbf{Testing:} pytest with FastAPI's
    \texttt{TestClient}.
\end{itemize}

\section{Implementation Details}

\subsection{Module 1: Identity and Shift Lifecycle}

The backend verifies each bearer token and resolves it to a
\texttt{staff} row, caching verified tokens briefly to avoid
a round trip per request.

Two authorisation questions are kept separate:
\texttt{ensure\_can\_view\_patient} (true also for a patient
being handed to her, since she must read the record before
deciding whether to accept) and
\texttt{ensure\_assigned\_patient} (true only while the
patient is hers). Collapsing them would either block the
receiving nurse from reading what she is being asked to
accept, or let her write to a patient who is not yet hers.

Shift handling is idempotent: \texttt{GET
/api/shifts/current} returns the open shift or opens one, so
a refresh cannot create duplicates. Closing a shift is
refused with HTTP 409 while patients remain assigned with no
handoff sent.

\subsection{Module 2: The Patient Record}

Each patient screen shows three views over the same rows:
\emph{Details}, \emph{Major Decisions}, and the full
chronological \emph{Log}. Notes accumulate through the
shift, which is what makes the handoff a review task rather
than a recall task.

Notes are append-only: no update or delete path exists in
the API or the RLS policies. A record that can be silently
edited afterwards cannot support the accountability the
system exists to provide.

\subsection{Module 3: Per-Patient Handoff}

The flow is \textbf{prepare}, \textbf{review},
\textbf{respond}.

Preparing gathers the shift's notes for each selected
patient and pre-fills a \emph{Details} and \emph{Major
Decisions} box by grouping them by section --- her own text,
saving retyping, writing nothing to the database. Reviewing
lets her edit either box and pick a named recipient. Sending
creates one \texttt{handoffs} row per patient with
\texttt{sent\_at} and a \texttt{sent} audit event.

The receiving nurse opens the handoff --- recording a
\texttt{viewed} event, which makes ``she read it before
acting'' auditable rather than assumed --- and reads
Details, Major Decisions and the Log before responding.

\diagramslot{[DIAGRAM SLOT --- Handoff state machine:
\texttt{pending} $\rightarrow$ \texttt{accepted} /
\texttt{declined} / \texttt{issue\_reported}, annotated with
whether the patient assignment transfers, and where
\texttt{responded\_at} is written.]}{Handoff State
Machine}{handoffstate}{5.5cm}

The three responses differ deliberately:

\begin{itemize}
  \item \textbf{Accept} transfers the patient.
  \item \textbf{Decline} requires a reason and leaves the
    patient with the sender --- wrong recipient, or too much
    missing to act on.
  \item \textbf{Report Issue} requires a reason and
    \emph{still transfers} the patient: something is wrong,
    but it needs resolving rather than rejecting, and the
    patient must not be left unattended meanwhile.
\end{itemize}

Every response writes \texttt{responded\_at} and appends an
event. The update is conditional on the row still being
pending, so a concurrent second response returns HTTP 409
rather than overwriting a recorded decision.

\subsection{Module 4: Frontend}

Three primary surfaces --- \emph{Shift}, \emph{My
Patients}, \emph{Handoff} --- with nested routes for a
patient, a note composer and a handoff review. Everything
except login sits behind a \texttt{RequireAuth} guard.

\diagramslot{[DIAGRAM SLOT --- Screen flow: Login
$\rightarrow$ Shift dashboard, My Patients $\rightarrow$
Patient detail $\rightarrow$ Note composer, Handoff
(Incoming / Send tabs) $\rightarrow$ Handoff review
$\rightarrow$ Accept / Decline / Report
Issue.]}{Application Screen Flow}{screenflow}{7cm}

The visual system is restrained: a single hue at varying
opacity, and semantic colour reserved for state that matters
--- amber pending, emerald accepted, rose declined. The
layout targets shared ward tablets, using responsive units
and safe-area insets.

% ===== CHAPTER 3: RESULTS & EVALUATION =====

\chapter{Results and Evaluation}

\section{Testing Methodology}

\begin{itemize}
  \item \textbf{Route tests.} A \texttt{pytest} suite drives
    the handoff state machine and note recording through
    FastAPI's \texttt{TestClient}, with identity injected by
    dependency override and the data layer replaced by
    recording doubles. Each test asserts on the writes
    actually attempted --- a status code cannot tell you
    whether declining left the patient with the sender.
  \item \textbf{Constraint tests.} Each database constraint
    was checked by attempting the forbidden operation
    against the live schema in a transaction and confirming
    rejection.
  \item \textbf{Integration checks.} Schema, foreign keys
    and seeded data were verified against the live Supabase
    instance, including signing in as a seeded nurse and
    confirming she resolves to the expected patients.
  \item \textbf{Build and security.} The frontend build
    type-checks the whole client before bundling; Supabase's
    advisor was run after each migration to catch missing
    RLS policies.
\end{itemize}

\section{Verification Results}

\begin{table}[h]
\centering
{\small
\begin{tabular}{|p{7.2cm}|r|l|}
  \hline
  \textbf{Check} & \textbf{Count} & \textbf{Result} \\
  \hline
  Automated backend tests & 15 & All pass \\
  \hline
  \quad Handoff state machine and access & 10 & All pass \\
  \hline
  \quad Note recording and access & 5 & All pass \\
  \hline
  Database constraint checks & 7 & All enforced \\
  \hline
  Tables with Row Level Security enabled & 7 of 7 & Enabled \\
  \hline
  Critical RLS security findings & 0 & Clear \\
  \hline
  Frontend type-check and build & 1 & Passes \\
  \hline
\end{tabular}}
\caption{Verification results}
\label{tab:verification}
\end{table}

The constraint checks confirmed, by direct attempt against
the live schema, that a self-addressed handoff is rejected;
a reasonless decline is rejected; a second pending handoff
for the same patient is rejected; resolving without a
response timestamp is rejected; accepting records the TTA
clock; and a patient becomes eligible for a fresh handoff
once the previous one resolves.

Access control was verified at both layers. At the database
layer, querying as a seeded nurse returned exactly her 3 of
the 6 seeded patients. At the application layer, tests
confirm a nurse cannot write notes for another nurse's
patient, cannot read a handoff addressed to someone else,
and that a sender cannot answer her own handoff on the
recipient's behalf.

RLS is enabled with policies on all seven tables and the
anonymous role's access revoked, so no table is reachable
--- or discoverable --- without a valid session.

The API surface is 16 endpoints across five routers (staff,
shifts, patients, notes, handoff).

\section{Evaluation Metric: Time-to-Acknowledgement}

The primary metric is \textbf{Time-to-Acknowledgement}: the
interval between a handoff being sent and answered. Every
handoff records \texttt{sent\_at} and \texttt{responded\_at},
and a constraint guarantees the pair stays consistent, so
the metric is computable directly:

\begin{verbatim}
SELECT avg(responded_at - sent_at)
FROM   handoffs
WHERE  responded_at IS NOT NULL;
\end{verbatim}

The secondary metric --- clean-accept rate, the proportion
accepted without an issue reported --- comes from the same
table.

\textbf{No measured TTA values are reported here.} That
requires a pilot with practising nurses, which has not been
run. Numbers from the development team exercising its own
interface would measure familiarity with that interface, not
handoff performance in a ward. The deliverable at this stage
is the instrumentation.

\section{Analysis}

All five objectives are implemented and verified.
Accountability is end-to-end: identity comes from a verified
token, every transition is recorded with its actor, and the
\texttt{viewed} event distinguishes a handoff that was read
from one merely accepted. Measurability is guaranteed by
constraint rather than convention. Access control is
enforced at two independent layers and demonstrated by tests
that assert on attempted writes.

\subsection{Challenges}
\label{sec:challenges}

\textbf{Concurrent database access.} Each screen issues
several API requests in parallel, and each request runs
several queries. The synchronous database client, driven
from a thread pool, shares one connection pool across those
threads; under concurrency this failed intermittently on
Windows, and because the failure took every in-flight
request on the connection with it, a whole screen would fail
at once. Moving to the asynchronous client removed the
shared-pool problem --- failures over a repeated
175-request load test went from 19 to zero.

\textbf{A permissions fault that looked like a data
problem.} With RLS enabled, backend reads failed with
\emph{permission denied}, which surfaced as a nurse simply
having no patients --- a symptom pointing away from its
cause. The backend had been configured with the publishable
key, which cannot bypass RLS. Besides using the correct key,
a startup check now refuses to launch on a publishable one,
so this cannot again present as an empty patient list.

\textbf{Defining ``Report Issue''.} Accept and Decline have
obvious semantics; this does not. Behaving like Decline
risks bouncing back a patient who is then unattended;
behaving like Accept risks the problem being forgotten. It
transfers care but keeps the issue on record, on the
reasoning that patient coverage must never be what lapses
while a documentation problem is sorted out.

\subsection{What Is Deferred}
\label{sec:deferred}

\begin{itemize}
  \item The \textbf{doctor interface} --- consolidated
    per-patient logs and nurse-to-doctor escalation. The
    \texttt{alerts} table and
    \texttt{assigned\_doctor\_id} already support it.
  \item The \textbf{administrator interface} --- staff
    management and assignment, which currently needs direct
    database access.
  \item \textbf{Voice capture} for dictating notes.
  \item \textbf{Push notification} of new or disputed
    handoffs; the inbox badge refreshes on navigation.
  \item \textbf{Deployment} --- the backend is not yet
    hosted.
  \item \textbf{The pilot}, and all measured outcome data.
\end{itemize}

% ===== CHAPTER 4: CONCLUSION =====

\chapter{Conclusion}

\section{Summary of Work}

CareBridge delivers a working nurse handoff workflow. A
nurse signs in against a real identity provider, sees the
patients under her care, and builds a structured record for
each across Details and Major Decisions. At shift end she
selects patients, reviews the text, and sends each to a
named colleague, who reads it in full and answers with
Accept, Decline or Report Issue --- with a reason where one
is owed. Every step is recorded with its actor and
timestamp.

All five objectives are met and verified by 15 automated
tests, 7 constraint checks, RLS across all seven tables and
a clean security advisory.

\section{Key Findings}

\begin{itemize}
  \item \textbf{Accountability needs an addressed artefact.}
    A handoff cannot be declined, disputed or measured until
    it is a durable object with a named sender and
    recipient. Making handoffs per-patient is what made the
    three responses coherent as a set.
  \item \textbf{Guarantees should be structural.} Rules in
    the database hold whatever code runs; rules in the
    interface hold only until a second path to the data
    exists.
  \item \textbf{Measurement must be designed in.} TTA is
    computable because two timestamps and a constraint were
    added at schema-design time; retrofitting a metric is
    far harder.
\end{itemize}

\section{Future Work and Recommendations}
\label{sec:futurework}

\begin{enumerate}
  \item \textbf{Complete the remaining actors} --- the
    doctor and administrator screens. Both are supported by
    the existing schema and need no migration.
  \item \textbf{Run the pilot and measure TTA} with 10--15
    nurses or a simulated ward, establishing a baseline for
    the current process by observation first so the
    comparison means something.
  \item \textbf{Add voice capture} for note entry, with a
    mandatory review step so no transcribed text reaches the
    record without the nurse approving it.
  \item \textbf{Deploy the backend} to managed hosting with
    CORS restricted to the deployed client origin.
  \item \textbf{Add push notification} so a nurse learns of
    a handoff without navigating to look for it.
\end{enumerate}

\section{Final Remarks}

The most useful lesson concerned where to put a rule. Rules
in the interface hold until someone builds a second path to
the same data; rules in the database hold unconditionally.
Working through this repeatedly changed the question we
asked from ``how do we make the app do the right thing'' to
``how do we make the wrong thing impossible''.

Handoffs will not stop being difficult because software
exists. But the specific failure this project targets ---
information that does not survive the transition between two
people --- is now recorded, addressed, answerable and
measurable.

% ===== BIBLIOGRAPHY =====

\printbibliography[heading=bibintoc]

\end{document}
